\section{Bradley-Terry model}
\label{sec:BTmodel}

Bradley–Terry models \cite{bradley_terry_1952} provide a natural starting point for modelling football match outcomes, as they directly relate the probability of winning to latent team skill parameters.
In their simplest form, these models treat matches as binary events, focusing only on the final outcome rather than the underlying scoreline.

In this work, we consider a family of Bradley–Terry–type models of increasing complexity, beginning with a basic binary formulation and progressively incorporating home effect and the possibility of draws.
These models serve both as a baseline for predictive performance and as a test for methodological validation, given their relatively low-dimensional parameter space and well-understood identifiability properties.

\subsection{General Formulation}
\label{sec:BT_general_formulation}

Before presenting the specific variants of the Bradley–Terry model, we introduce a general formulation that incorporates all its structural components: team skills, home effect, and the possibility of draws.
This serves as a conceptual starting point, from which different models naturally emerge through parameter restrictions and simplifying assumptions.

The Bradley-Terry model can be expressed in terms of team strengths and outcome probabilities.
For a match between team $i$ (home) and team $j$ (away), the strengths are defined as:
\begin{align}
    \label{eq:bt_general_home_strength}
    s_i &= \exp(\theta_i + \nu) \\
    \label{eq:bt_general_away_strength}
    s_j &= \exp(\theta_j),
\end{align}

\noindent where:
\begin{itemize}
    \item $\theta_i \in \mathbb{R}$ is the log-skill parameter for team $i$;
    \item $\nu \in \mathbb{R}$ is the home effect parameter;
    \item $s_i$ is the strength of team $i$ when playing at home;
    \item $s_j$ is the strength of team $j$ when playing away.
\end{itemize}

The probability of each outcome is then given by:
\begin{align}
    \label{eq:bt_general_prob_home}
    P(\text{team } i \text{ wins}) &= \frac{s_i}{s_i + s_j + s_{\text{tie}}} \\
    \label{eq:bt_general_prob_away}
    P(\text{team } j \text{ wins}) &= \frac{s_j}{s_i + s_j + s_{\text{tie}}} \\
    \label{eq:bt_general_prob_tie}
    P(\text{draw}) &= \frac{s_{\text{tie}}}{s_i + s_j + s_{\text{tie}}},
\end{align}

\noindent where $s_{\text{tie}} = \kappa \sqrt{s_i \cdot s_j}$ is the tie strength, with $\kappa \geq 0$ being the tie parameter, following the approach of \textcite{davidson_1970}.
When $\kappa = 0$, draws are not allowed and the model reduces to a binary outcome model, with the outcome $y_{i,j}$ following a Bernoulli distribution.
When $\kappa > 0$, the outcome $y_{i,j}$ follows a categorical distribution with these probabilities.

Given a set of matches $\mathcal{M}$, the likelihood can be computed as:
\begin{equation}
    \label{eq:bt_general_likelihood}
    \mathcal{L}(\boldsymbol{\theta}, \nu, \kappa ; \mathbf{y}) = \prod_{m \in \mathcal{M}} P(y_m \mid \theta_{i_m}, \theta_{j_m}, \nu, \kappa),
\end{equation}

\noindent where $y_m$ is the outcome of match $m$ between teams $i_m$ and $j_m$.

The Bradley-Terry model defined by equations \ref{eq:bt_general_prob_home}, \ref{eq:bt_general_prob_away}, and \ref{eq:bt_general_prob_tie} presents identifiability issues, with the likelihood remaining invariant under translations of all log-skill parameters by a constant $c$:
\begin{equation}
    \mathcal{L}(\boldsymbol{\theta}, \nu, \kappa ; \mathbf{y}) = \mathcal{L}(\boldsymbol{\theta} + c, \nu, \kappa ;\mathbf{y}),
\end{equation}

\noindent since:
\begin{align*}
    (\theta_i + c) + \nu - (\theta_j + c) &= \theta_i + \nu - \theta_j, \\
    \exp(\theta_i + c + \nu) &= \exp(c) \cdot \exp(\theta_i + \nu), \\
    \exp(\theta_j + c) &= \exp(c) \cdot \exp(\theta_j),
\end{align*}

\noindent and the constant $\exp(c)$ cancels out in the ratio of strengths.

To ensure model identifiability, we impose a sum-to-zero constraint on the log-skill parameters:
\begin{equation}
    \label{eq:sum_zero_cons_bt}
    \sum_{i=1}^{n} \theta_i = 0.
\end{equation}

This constraint is implemented by parametrizing the vector as:
\begin{equation}
    \boldsymbol{\theta} = \lrc{\theta_1, \theta_2, \ldots, \theta_{n-1}, -\sum_{i=1}^{n-1} \theta_i},
\end{equation}

\noindent where only the first $n-1$ parameters are estimated, and the last is determined by the constraint.

\subsection{Model Variants}
\label{sec:BT_model_variants}

Different models arise from simplifying assumptions about the parameters in the general formulation.
The models can be categorized based on two key dimensions:
\begin{enumerate}
    \item \textbf{Home Effect}: Present ($\nu \neq 0$) vs. Absent ($\nu = 0$);
    \item \textbf{Ties Allowed}: Present ($\kappa > 0$) vs. Absent ($\kappa = 0$).
\end{enumerate}

This section describes the four Bradley-Terry model variants (BT1--BT4) used to model match outcomes in football matches in the present work.
These models are organized from the simplest binary outcome model without home effect, BT1, to the most complex ternary outcome model with home effect, BT4.
The model names were chosen for simplicity, avoiding unnecessary repetition.

Model \textbf{BT1} is the simplest Bradley-Terry model, using a single log-skill parameter $\theta_i$ per team with no home effect and no ties allowed.
The log-odds of team $i$ winning against team $j$ is given by $\log(\text{odds}) = \theta_i - \theta_j$, which corresponds to a binary outcome model where $P(\text{team } i \text{ wins}) = \exp(\theta_i) / (\exp(\theta_i) + \exp(\theta_j))$.

Model \textbf{BT2} extends BT1 by incorporating a home effect parameter $\nu$.
The log-odds becomes $\log(\text{odds}) = \theta_i + \nu - \theta_j$ for the home team, capturing the additional effect of playing at home.
Like BT1, this model does not allow ties.

Model \textbf{BT3} extends BT1 by allowing ties through the introduction of a tie parameter $\kappa \geq 0$.
The tie strength is defined as $s_{\text{tie}} = \kappa \sqrt{s_i \cdot s_j}$, where $s_i = \exp(\theta_i)$ and $s_j = \exp(\theta_j)$.
This model does not include home effect, so the strengths are symmetric with respect to home and away contexts.

Model \textbf{BT4} combines the features of BT2 and BT3, incorporating both home effect and ties.
The home team's strength becomes $s_i = \exp(\theta_i + \nu)$, while the away team's strength remains $s_j = \exp(\theta_j)$.
The tie strength is $s_{\text{tie}} = \kappa \sqrt{s_i \cdot s_j}$, providing the most flexible specification that accounts for both home effect and the possibility of draws.

The taxonomy of these model variants is illustrated in Figure~\ref{fig:bt_models_taxonomy}.

\begin{figure}[ht]
    \centering
    \begin{tikzpicture}[
        level 1/.style = {sibling distance = 4cm, level distance = 2.5cm},
        level 2/.style = {sibling distance = 2cm, level distance = 2cm},
        every node/.style = {
            draw,
            rectangle,
            rounded corners,
            align=center,
            minimum width = 1.5cm,
            minimum height = 0.8cm,
            font = \scriptsize
        },
        root/.style = {
            minimum width = 3.5cm,
            minimum height = 1.5cm,
            font = \normalsize
        },
        leaf/.style = {
            fill = gray!20,
            font = \tiny
        }
    ]
        \node[root] {
            $\begin{aligned}
                s_i &= \exp(\theta_i + \nu) \\
                s_j &= \exp(\theta_j) \\
                s_{\text{tie}} &= \kappa \sqrt{s_i \cdot s_j}
            \end{aligned}$ \\
            General Formulation
        }
            child {
                node {
                    $\kappa = 0$ \\
                    No Ties
                }
                child {
                    node {
                        $\nu = 0$ \\
                        No Home \\
                        Effect
                    }
                    child {
                        node[leaf] {BT1}
                    }
                }
                child {
                    node {
                        $\nu \neq 0$ \\
                        With Home \\
                        Effect
                    }
                    child {
                        node[leaf] {BT2}
                    }
                }
            }
            child {
                node {
                    $\kappa > 0$ \\
                    Ties Allowed
                }
                child {
                    node {
                        $\nu = 0$ \\
                        No Home \\
                        Effect
                    }
                    child {
                        node[leaf] {BT3}
                    }
                }
                child {
                    node {
                        $\nu \neq 0$ \\
                        With Home \\
                        Effect
                    }
                    child {
                        node[leaf] {BT4}
                    }
                }
            };
    \end{tikzpicture}
    \caption{
        \textbf{Taxonomy of the four Bradley-Terry model variants}.
        The root node displays the complete specification with log-skill parameters ($\theta_i$), home effect ($\nu$), and tie parameter ($\kappa$).
        The first level of branching corresponds to whether ties are allowed ($\kappa > 0$) or not ($\kappa = 0$).
        The second level distinguishes models with ($\nu \neq 0$) and without ($\nu = 0$) home effect.
        Each leaf node represents a specific model variant, with the model identifier (BT1--BT4) indicating the final configuration after all simplifications have been applied.
    }
    \label{fig:bt_models_taxonomy}
\end{figure}
